% FaST-ISL paper draft (ACL format). Build: cd paper && tectonic main.tex
% Everything marked \todo{} needs the real iSign experiments (ROADMAP.md, step 10).
\documentclass[11pt]{article}
\usepackage[review]{acl}
\usepackage{times}
\usepackage{latexsym}
\usepackage[T1]{fontenc}
\usepackage[utf8]{inputenc}
\usepackage{microtype}
\usepackage{booktabs}
\usepackage{amsmath}
\usepackage{amssymb}
\usepackage{graphicx}
\usepackage{xcolor}
\newcommand{\todo}[1]{\textcolor{red}{[TODO: #1]}}
\newcommand{\fast}{FaST-ISL}

\title{\fast: Face-aware, Speed-robust, Streaming Translation of\\Continuous Indian Sign Language}

\author{Anonymous submission}

\begin{document}
\maketitle

\begin{abstract}
Indian Sign Language (ISL) marks questions, negation and emphasis with the face and head, yet every
recent continuous ISL-to-English translation system we survey (including all ISL systems at WSLP 2026)
uses body and hand keypoints only. These systems are also trained and tested at one natural signing
speed, and they translate pre-cut clips rather than live, unsegmented signing. We present \fast{}, a
pose-to-text translator that adds (i)~a gated face and head-pose stream, (ii)~speed robustness via a
multi-rate temporal encoder and 0.5--2$\times$ temporal resampling, and (iii)~a learned, causal
sentence-boundary detector that turns continuous signing into sentence-level English text and speech in
real time, even when the signer is fast or does not pause. A pause-threshold baseline, by contrast,
merges sentences as soon as signing speeds up. On
the iSign benchmark, \fast{} improves BLEU-4 from 1.80 to 2.07 (+15\% relative, $p<0.01$) over a
body+hands baseline with the same decoder, keeps 95\% of its BLEU-4 when signing is twice as
fast (baseline: 85\%), and recovers more of the question and negation forms (recall 35\% vs.\ 27\% and
55\% vs.\ 42\%). A region
ablation shows that \todo{region} carries most of the benefit. The full system runs at \todo{fps} on
a laptop CPU with \todo{ms}~ms model latency per sentence. Code, configs and the demo are released.
\end{abstract}

\section{Introduction}

A person signing to a camera does not stop between sentences, does not sign at a fixed speed, and
does not express everything with the hands. Sign languages use \emph{non-manual markers}---eyebrow
raises, head shakes, mouth shapes---as grammar: in ISL, raised eyebrows typically mark a yes/no question
and a head shake marks negation \citep{joshi2024isign}. A translator that discards the face cannot, in
principle, distinguish \emph{``You are going home.''} from \emph{``Are you going home?''} when the hands
are identical.

The release of iSign \citep{joshi2024isign}, with more than 118k ISL video--English sentence pairs, and
the WSLP shared tasks \citep{wslp2026findings} made continuous ISL-to-English translation a benchmarked
task. Quality is still low: the best reported iSign BLEU-4 is below 10 \citep{roychowdhury2025motion},
and the 2026 shared-task overview notes that official BLEU can sit at zero.

We reviewed the recent ISL recognition and translation literature, including all 24 papers of the WSLP
2026 proceedings, and find three gaps (\S\ref{sec:related}):
\begin{enumerate}
\item \textbf{The face is missing.} Pose-based ISL translators keep 33 body and 42 hand points and drop
  the face \citep{dhamecha2026temporal, ahmed2026islim, shrestha2026contrastive, geetha2025signflow}.
  Face-aware translation has helped Chinese and German sign language \citep{tu2026feaslt} but has not
  been tested on ISL.
\item \textbf{Speed is fixed.} No ISL translation work measures robustness to faster signing, although
  signing rate varies across signers and situations.
\item \textbf{No streaming.} ISL systems translate pre-segmented clips; real-time systems are
  closed-domain and output glosses \citep{geetha2025signflow} or take seconds per response
  \citep{nedungadi2025islca}.
\end{enumerate}

\paragraph{Contributions.}
\begin{itemize}
\item \fast{}, a pose-to-English translator with a \emph{gated} face/head stream that can open only when
  the face carries meaning, and an ablation of which facial region (mouth, eyes, eyebrows, head pose)
  matters for ISL translation.
\item A speed-robust encoder (multi-rate temporal convolution plus 0.5--2$\times$ resampling) and a
  \emph{speed-perturbed} iSign test protocol (1.5$\times$, 2$\times$).
\item A non-manual evaluation: question- and negation-form accuracy on iSign, alongside BLEU, chrF and
  ROUGE-L.
\item A streaming system: a learned causal sentence-boundary detector trained on stitched streams with
  random speed and pause length, a stream-level evaluation on reassembled iSign videos, and an
  open-source web demo producing English text and speech. We show that a pause-threshold segmenter
  fails on fast or pause-free signing while the learned detector does not.
\end{itemize}

\section{Related Work}\label{sec:related}

\paragraph{ISL translation.} iSign \citep{joshi2024isign} provides MediaPipe Holistic poses and
English sentences; its baselines reach BLEU-4 around 1.5. PoseStitch-SLT \citep{posestitch2025} and
motion-aware modelling \citep{roychowdhury2025motion} raise this to 3.4 and 8.6. At WSLP 2026,
\citet{dhamecha2026temporal} feed body and hand keypoints with frame differences into T5;
ISLiM \citep{ahmed2026islim} is a compact end-to-end model (iSign BLEU-4 5.85);
\citet{shrestha2026contrastive} add contrastive grounding; \citet{elaswad2026articulator} transfer ASL
models to ISL. None of these uses face features, varies signing speed, or translates a live stream.

\paragraph{Real-time ISL systems.} SignFlow \citep{geetha2025signflow} recognises continuous ISL in a
fixed e-governance vocabulary from RGB, pose and hand input and outputs glosses. ISL-CA
\citep{nedungadi2025islca} is a closed-domain chatbot with 6.5\,s mean response time.
\citet{srivastava2024continuous} classify a small set of phrases from Holistic landmarks including the
face, but do not translate open sentences.

\paragraph{Non-manual features.} \citet{pham2025facial} find the mouth to be the most informative facial
region for recognition and report that few prior works used the face. FEA-SLT \citep{tu2026feaslt}
adds facial expressions to gloss-free translation on PHOENIX14T and CSL-Daily. A review of low-resource
sign language processing lists the neglect of non-manual features among its core challenges
\citep{lowresource2026review}. Our work brings this question to ISL translation from pose and asks
which region helps.

\paragraph{Speed perturbation.} Changing playback speed is a standard augmentation in speech recognition
\citep{ko2015audio}. We apply it to sign keypoints and, unlike prior work, also use it to build
fast-signing \emph{test} sets.

\section{Method}

\subsection{Input features}
Each frame holds MediaPipe Holistic \citep{lugaresi2019mediapipe} keypoints: 33 body points, 21 points
per hand, and a 152-point face subset (lips 40, eyes 32, eyebrows 20, face outline 36, nose 24) chosen
from the 468-point mesh. We build five streams per frame $t$:
\begin{itemize}
\item \textbf{body}: 25 upper-body points, centred on the shoulder midpoint and divided by shoulder width;
\item \textbf{hands}: each hand in body coordinates and relative to its own wrist (scaled by palm size);
\item \textbf{face}: face points relative to the face centre, divided by the outer-eye distance, so the
  stream encodes expression independently of head position;
\item \textbf{head}: six head-pose proxies: face centre and scale relative to the shoulders, roll (eye
  line angle), and yaw/pitch proxies (nose offset from the face centre);
\item \textbf{present}: four flags marking whether body, each hand and face were detected.
\end{itemize}
Undetected parts are zero and their normaliser is carried over from the last valid frame. Each stream
$x$ is fed as $[x_t, x_t - x_{t-1}]$ so motion is explicit.

\subsection{Gated face fusion}
Manual and non-manual inputs are projected separately (LayerNorm--Linear--GELU--Linear):
$m_t = P_b(\text{body}_t) + P_h(\text{hands}_t) + W_p\,\text{present}_t$ and
$f_t = P_f(\text{face}_t) + P_{hp}(\text{head}_t)$. Head pose has its own projection: its six numbers
would otherwise be swamped by the 304 face coordinates. A per-channel gate decides how much face
information enters:
\begin{equation}
  g_t = \sigma\!\left(W_g [m_t; f_t] + b_g\right), \qquad h_t = m_t + g_t \odot f_t .
\end{equation}
$b_g$ is initialised to $-1$ so the gate starts mostly closed and opens only where the face helps. The
ablation \emph{concat} replaces the gate by $h_t = W[m_t; f_t]$; the \emph{baseline} uses $h_t = m_t$.

\subsection{Speed-robust temporal encoder}
A multi-rate block applies three depthwise temporal convolutions (kernel 5, dilations 1, 2, 4) in
parallel, mixes them with a $1{\times}1$ convolution and a residual connection, and halves the frame
rate with a stride-2 convolution. The same movement is thus seen at three time scales, so a sign made
twice as fast still matches a learned pattern. Four pre-norm Transformer layers ($d{=}256$, 4 heads)
follow. During training each clip is, with probability 0.5, resampled to a speed factor drawn
log-uniformly from $[0.5, 2]$ by linear interpolation between frames.

\subsection{Decoder}
A linear bridge maps encoder states to the input embedding space of a pretrained T5-small
\citep{raffel2020t5}, whose encoder and decoder are fine-tuned end to end, following
\citet{dhamecha2026temporal}. We decode with beam size 4 and block repeated 3-grams.

\paragraph{Keeping the decoder grounded.} With a pretrained decoder and only $\sim$10 words per target,
our first baseline drifted to fluent, topic-level guesses (\emph{``Page number 76''} for a textbook
clip; generic news for news clips) and scored no worse on 1.5$\times$ replays than on normal speed, a
sign that it barely read the input. Two training changes, used for every system in
Table~\ref{tab:main}, fix this. \emph{Decoder word dropout}: 25\% of the decoder's input tokens are
replaced by \texttt{<unk>} during training, so the next word cannot be predicted from the language prior
alone. \emph{Clip--sentence contrastive loss}: a symmetric InfoNCE term (weight 0.3) between the
mean-pooled sign encoding and the T5 encoding of the target sentence, where identical sentences in a
batch count as positives. Together they raise the baseline from 0.84 to 1.80 BLEU-4.

\subsection{Streaming segmentation}
\paragraph{Pause rule (baseline).} For live use, a segmenter receives keypoints at 10\,fps. It tracks the hands' motion energy
$e_t = \max_{\text{hand}} \lVert \bar{p}_t - \bar{p}_{t-1} \rVert / w_t$ (mean hand position, shoulder
width $w_t$), smoothed by an exponential moving average. A sentence starts when a hand is in the
signing space and $e_t$ exceeds $\tau_{on}$, and ends after 0.6\,s with $e_t < \tau_{off}$ (hysteresis)
or with the hands lowered. Sentences longer than 12\,s are cut at the calmest recent frame to bound
latency. Each finished segment is translated at once and spoken with text-to-speech.

\paragraph{Why a pause rule is not enough.} When a person signs twice as fast, a 0.8\,s pause lasts
0.4\,s and falls under any fixed threshold, so consecutive sentences merge into one long segment that
the translator cannot handle (Table~\ref{tab:synthetic-stream}). Fluent signers also often move
straight into the next sentence without pausing.

\paragraph{Learned boundary detector.} We therefore train a small causal network (0.4M parameters) to
predict sentence ends frame by frame. It reads the body, hand, head-pose and presence streams (with
velocities), applies a causal temporal convolution (kernel 5) and a two-layer GRU ($d{=}128$), and
outputs $p_t$, the probability that a sentence ended at frame $t-\delta$. The look-ahead $\delta = 3$
frames (0.3\,s) lets it see the start of the next sentence. Training streams join 2--5 consecutive
training sentences of the same video; each stream is replayed at a speed drawn log-uniformly from
$[0.5, 2]$ and each gap gets a pause of 0--1.2\,s, or none at all with probability 0.35. Targets are 1
within $\pm1$ frame of each sentence end; we use binary cross-entropy with positive weight 8. At run time
the GRU state is carried across frames, a boundary is emitted at local maxima of $p_t$ above 0.5 (with a
0.6\,s refractory period), and still frames at both ends of a segment are trimmed so pause holds are not
translated.

\section{Experimental Setup}

\paragraph{Data.} iSign v1.1 \citep{joshi2024isign}: 127{,}237 clips with poses and a non-empty
sentence; the test split has 5{,}325 clips. iSign
clip identifiers have the form \texttt{<video>-<n>}; we hash the source video to assign 90/5/5\% of
\emph{videos} to train/validation/test, so no test sentence comes from a video seen in training.
Frames are kept at 10\,fps and at most 256 per clip.

\paragraph{Systems.} All systems share the decoder, optimiser and data (Table~\ref{tab:main}).
\textbf{B0}: body + hands, single-rate encoder, no speed augmentation (the setting of current ISL
work). \textbf{B1}: B0 + multi-rate encoder + speed augmentation. \textbf{F-*}: B1 + one facial region
(mouth / eyes / eyebrows / head pose = outline + nose). \textbf{\fast}: all regions, gated.
\textbf{\fast-concat} and \textbf{\fast{} w/o aug} isolate the gate and the augmentation.

\paragraph{Training.} AdamW (encoder lr $3{\cdot}10^{-4}$, T5 lr $10^{-4}$, weight decay 0.01), 1k
warm-up steps then cosine decay, batch 32, label smoothing 0.1, 20 epochs, keeping the checkpoint with
the best validation BLEU-4. Each model trains in fp32 on one NVIDIA T4 (16\,GB), about 25 minutes per
epoch; two models train side by side on a 2$\times$T4 machine.

\paragraph{Metrics.} BLEU-4 \citep{papineni2002bleu} with SacreBLEU \citep{post2018sacrebleu}, chrF
\citep{popovic2015chrf} and ROUGE-L \citep{lin2004rouge}. \emph{Speed robustness}: the test set
replayed 1.5$\times$ and 2$\times$ faster. \emph{Non-manual accuracy}: a reference is a question if it
ends with ``?'' and is negated if it contains a negation word (\emph{not, no, never, n't, \ldots});
question (negation) accuracy is the share of test sentences where the output agrees with the reference
on this property. These automatic tags are a proxy; \S\ref{sec:limits} discusses their limits.
\emph{Streaming}: we rebuild continuous signing by joining five consecutive test sentences of the same
video with 0.8\,s, 0.4\,s or no pause, replay each stream at 1$\times$, 1.5$\times$ and 2$\times$, and
report sentence-boundary F1 (tolerance 0.5\,s), stream-level BLEU-4 and per-sentence latency for the pause
rule and the learned detector.

\section{Results}

% Placeholder: overwritten by `python -m islt.report runs --out paper/results_table.tex`
\begin{table*}[t]\centering\small
\begin{tabular}{lrrrrrrrr}\toprule
& \multicolumn{5}{c}{iSign test} & & \multicolumn{2}{c}{Non-manual acc.} \\
Model & BLEU-4 & $\Delta$ & chrF & ROUGE-L & B4@1.5$\times$ & B4@2$\times$ & Question & Negation \\
\midrule
\multicolumn{9}{c}{\textcolor{red}{[TODO: run scripts/run\_ablations.sh or colab/train\_fastisl.ipynb]}} \\
\bottomrule\end{tabular}
\caption{Main results and ablations on the iSign test split (signer-video disjoint).}
\label{tab:main}\end{table*}


\paragraph{Main result.} \fast{} reaches 2.07 BLEU-4 against 1.80 for the body+hands baseline, a 15\%
relative gain, and also improves chrF (17.20 vs.\ 16.90) and ROUGE-L (14.26 vs.\ 13.64). A paired
bootstrap over the 5{,}325 test clips (1{,}000 resamples) gives a BLEU-4 gain of $+0.27$ (95\% CI
$[0.11, 0.43]$, $p=0.001$) and a chrF gain of $+0.30$ ($[0.13, 0.46]$, $p<0.001$). Both systems exceed
the 1.47 BLEU-4 of the original iSign pose baseline \citep{joshi2024isign}, but remain far below the best
shared-task systems, so these are gains over a matched baseline, not a new state of the art. \fast{}
differs from B0 in both the face stream and the speed-robust encoder;
\todo{B1 (speed only) separates the two: FaST vs B1 = face, B1 vs B0 = speed robustness.}

\paragraph{Which facial region matters?} \todo{Compare F-mouth, F-eyes, F-eyebrows, F-head; relate to
\citet{pham2025facial} (mouth most informative for recognition).}

\paragraph{Fast signing.} At 2$\times$ speed the baseline keeps 85\% of its BLEU-4 (1.53) while
\fast{} keeps 95\% (1.96); at 1.5$\times$ neither loses accuracy. The gap at 2$\times$ is larger than at
normal speed: $+0.43$ BLEU-4 (95\% CI $[0.27, 0.59]$) and $+2.22$ chrF ($[2.03, 2.41]$), both $p<0.001$.

\paragraph{Non-manual grammar.} Question and negation \emph{accuracy} barely moves (94.3$\to$94.8\%
and 85.8$\to$86.3\%) because most test sentences are neither, so always predicting a statement already
scores high. \emph{Recall} is more telling: \fast{} produces a question for 35\% of the 169 question
references (baseline 27\%) and a negation for 55\% of the 530 negated references (baseline 42\%), and
BLEU-4 on the question subset rises from 0.85 to 1.19. \todo{Confirm against B1, and add examples where
FaST gets the question form right and B0 does not.}

\paragraph{Gate behaviour.} \todo{Mean gate value per frame; is it higher on question/negation
sentences?}

\paragraph{Streaming.} \todo{iSign version of Table~\ref{tab:synthetic-stream}: boundary F1, stream
BLEU-4 and latency for each pause length and speed, pause rule vs.\ learned detector; FPS of the full demo
on a laptop CPU (MediaPipe + model: 26\,ms per frame measured so far on an Apple M3).}

\paragraph{Sanity check on controlled data.} Before training on iSign we verified the pipeline on a
synthetic sign language in the same keypoint format, where some meaning is carried \emph{only} by the
face: raised eyebrows mark a question and a head shake inserts ``not''. Appendix~\ref{app:synthetic}
shows that the face stream and speed augmentation behave as intended there, and that on continuous
streams the learned boundary detector keeps boundary F1 at 100 while the pause rule drops to 40 at
2$\times$ speed. These numbers test the implementation, not ISL.

\section{Limitations}\label{sec:limits}
Face keypoints from MediaPipe are noisy under occlusion by the hands, which are often in front of the
face in ISL. Our non-manual tags come from English punctuation and negation words, not from annotation of
the signing, so they measure whether the output has the right sentence type, not whether the model used
the right cue. A hand-annotated non-manual test subset with ISL interpreters is future work. Streams are
built by joining clips cut from the same video, which approximates but does not equal natural continuous
signing. iSign is released for research use only; a deployed product needs separately consented data.
Our evaluation does not include Deaf users; a user study is planned.

\section*{Ethics Statement}
Sign language technology should be built with Deaf communities. Our demo processes video locally and
sends only keypoints to the model. Errors in translation could mislead in important settings (health,
law), so the system is a research prototype, not a replacement for interpreters.

\bibliography{references}

\appendix
\section{Synthetic sanity check}\label{app:synthetic}
The toy language has 20 words; each word is a distinct right-hand path and hand shape. Sentences have
3--4 words; 30\% are questions (eyebrows raised, text ends with ``?'') and 30\% negations (head shakes,
``not'' inserted after the subject, with no manual sign). Training/validation/test: 2000/200/300
sentences. All models use the same small configuration and 12 epochs.

% Filled by: python -m islt.report runs --synthetic --out paper/synthetic_table.tex
\textcolor{red}{[TODO: synthetic table]}


Table~\ref{tab:synthetic} isolates each component. Speed augmentation (B0$\to$B1) raises BLEU-4 at
2$\times$ speed from 16.1 to 53.7 without changing normal-speed quality. Adding the face and head
streams (B1$\to$FaST) lifts question and negation accuracy from chance level (68--69\%, i.e.\ always
predicting the majority form) to 100\%. Without speed augmentation the face-aware model is perfect at
normal speed but falls to 28.1 BLEU-4 at 2$\times$. In an earlier version, the six head-pose values
shared one projection with the 304 face coordinates and negation stayed at chance; giving head pose its
own projection fixed this, which motivated the design in \S3.2.

For streaming (Table~\ref{tab:synthetic-stream}), test sentences are joined four at a time. The pause
rule only works with long pauses at normal speed; the learned detector finds every boundary in all six
settings, and the remaining drop in stream BLEU-4 at 2$\times$ comes from the translator, not from
segmentation. Mean model latency was 66--100\,ms per sentence on an Apple M3.

\begin{table}[h]\centering\small
\begin{tabular}{llrrrr}\toprule
& & \multicolumn{2}{c}{Boundary F1} & \multicolumn{2}{c}{Stream BLEU-4} \\
Pause & Speed & rule & learned & rule & learned \\ \midrule
0.8\,s & 1$\times$ & 90.0 & 100.0 & 99.1 & 98.8 \\
0.8\,s & 2$\times$ & 40.0 & 100.0 & 2.6 & 87.4 \\
0.4\,s & 1$\times$ & 51.5 & 100.0 & 33.6 & 98.8 \\
0.4\,s & 2$\times$ & 40.0 & 100.0 & 3.1 & 87.7 \\
0.0\,s & 1$\times$ & 66.0 & 100.0 & 26.6 & 98.8 \\
0.0\,s & 2$\times$ & 40.0 & 100.0 & 2.4 & 86.6 \\
\bottomrule\end{tabular}
\caption{Continuous signing (synthetic, streams of 4 sentences). Boundary F1 and stream BLEU-4 for the pause rule vs.\ the learned boundary detector.}
\label{tab:synthetic-stream}\end{table}


\end{document}
